% RESEARCH_DRAFT_SHA256: 0de9ceba8a8cec8fdcbc6073ce241907c506404c923ce1d6a9c685a62602a7be
% Layout and source-reference wrappers expanded for the web renderer.
\documentclass[11pt]{article}
\usepackage[margin=25mm]{geometry}
\usepackage{amsthm}
\usepackage{fontspec,unicode-math}
\setmainfont{Latin Modern Roman}
\setmathfont{Latin Modern Math}
\usepackage{amsmath,amssymb,mathtools}
\usepackage{xurl,hyperref}
\hypersetup{colorlinks=true,urlcolor=blue,linkcolor=blue}
\setcounter{secnumdepth}{0}
\setlength{\parindent}{0pt}
\setlength{\parskip}{5pt}
\setlength{\emergencystretch}{3em}
\newcommand{\sref}[2]{\hyperlink{source:#1:#2}{[S#2]}}
\newcommand{\cataloguescope}[1]{\paragraph{Source question.}#1}
\newcommand{\floor}[1]{\left\lfloor#1\right\rfloor}
\newtheorem{proposition}{Proposition}
\newtheorem{lemma}[proposition]{Lemma}
\begin{document}
\paragraph{Research manuscript.} OpenAI \texttt{gpt-6-astra}, reasoning \texttt{ultra}; 11 October 2026.
\paragraph{Preserved source body.} The complete catalogue mathematical text and sources follow with layout and source links expanded, including source qualifications and its historical preparation notes. New research begins at the section headed Research attempt.

\section{Characterization of Ideal $0$–$1$ Matrices}\label{statement}

\subsection{Definitions and mathematical statement}

\paragraph{Shared notation.}Standard number systems are denoted by $\mathbb N,\mathbb Z,\mathbb Q,\mathbb R,\mathbb C$; index ranges, computational models and specialized notation follow the cited problem source and the definitions below.


A polyhedron is integral when it is the convex hull of its integer points; for a pointed polyhedron this is equivalent to all vertices being integral. $\mathbf1$ is the all-ones vector. Minimally nonideal refers to nonideality that disappears under the source\textquotesingle s prescribed minor operations.
\paragraph{Statement recorded in the supplied source.}
Characterize ideal $0$–$1$ matrices $M$, meaning those for which $\{x:Mx\geq\mathbf1,\ x\geq0\}$ is integral; in particular, seek an explicit characterization of minimally nonideal matrices. \href{https://doi.org/10.4171/owr/2005/50}{[S2]}
\paragraph{Source status.}
The supplied catalogue labels this record open. Its source-status metadata was checked against the supplied snapshot on 2026-10-11. This import does not independently establish current open status or mathematical correctness. The catalogue literature assessment is dated 2026-08-21.
\subsection{Short English statement}
Describe exactly which binary covering matrices define integral polyhedra and which are minimal obstructions.
\subsection{Sources}
\begin{itemize}
\item[S1] Ulam AI and the UnsolvedMath Contributors, UnsolvedMath: A Curated Collection of Open Mathematics Problems, catalogue entry. Catalogue curation and metadata: CC BY 4.0. The supplied source question is retained; definitions and the English summary are editorial additions.
\url{https://huggingface.co/datasets/ulamai/UnsolvedMath}.
\item[S2] Combinatorial Optimization (2006), Oberwolfach Reports, DOI 10.4171/owr/2005/50, pp. 2833--2932
\url{https://doi.org/10.4171/owr/2005/50}.
\textit{Underlying Oberwolfach Reports material retains its source-specific licence, including CC BY-SA 4.0 where applicable; see THIRD\_PARTY\_NOTICES.md.}
\end{itemize}
\paragraph{Preparation and review.}
Original question attributed to the cited source and UnsolvedMath contributors. Codex (GPT-6; exact serving revision not exposed) assisted with definitions, English summaries, source inspection and formatting. No new human mathematical authorship or independent proof verification is asserted. Catalogue identity checked against the complete supplied registry and existing repository statements; known flagged repetitions and obvious identical conjectures excluded. Full mathematical-equivalence/literature review has not been performed. Statement preparation and final review followed contribution prompts C and D; automated checks verify structure and identity, not open status.
\label{end:statement}

\section{Research attempt}
\subsection{Literature, source context and scope}
The full original catalogue body and source qualifications are preserved above. The relevant primary workshop passages were read in the cached publisher report. Historical claims and subsequent attributed results are distinguished from the mathematical work below. A restricted example, standard lemma or known theorem is not counted as new progress on the entire original target. Missing assumptions and repeated or rhetorical extracts are not silently turned into different solved problems.

\subsection{The row-size-two subclass and failure of a total-unimodularity shortcut}
For a simple graph $G$, let $M$ be its edge-versus-vertex incidence matrix, with two ones in every row. Then the source's covering polyhedron is
\[
P(G)=\{x\ge0:x_u+x_v\ge1\text{ for every edge }uv\}.
\]
We examine the restricted row-size-two case before attempting a general characterization of ideal matrices.
\begin{proposition}
Every vertex of $P(G)$ is integral if and only if $G$ is bipartite.
\end{proposition}
\begin{proof}
First suppose $G$ is bipartite. At a vertex, no coordinate can exceed 1: all incident inequalities would then be strict, allowing a sufficiently small positive and negative change of that coordinate while preserving feasibility. Let $F$ be the coordinates strictly between 0 and 1. A tight edge meeting $F$ has both endpoints in $F$. In any connected component of the graph of tight edges on $F$, assign opposite perturbations $+\epsilon,-\epsilon$ to the two bipartition classes, leaving all other coordinates fixed. Every tight edge remains tight, including isolated vertices of this tight-edge graph, which have no such condition. Finitely many nontight inequalities and positive coordinates remain feasible for both signs when $\epsilon>0$ is sufficiently small. If $F$ were nonempty this would express $x$ as the midpoint of two distinct feasible vectors, contradicting its being a vertex. Hence every coordinate is 0 or 1.
Conversely, a nonbipartite graph has an odd cycle of length $2r+1$. Put $x=1/2$ on its vertices and $x=1$ elsewhere; this is feasible, even if the cycle has chords. Give cycle coordinates cost 1 and all other coordinates a positive cost $\delta$ chosen so that their total cost is less than $1/2$. The displayed fractional vector has cost less than $r+1$. Every integral feasible vector must have positive value on a vertex cover of the cycle. Summing its $2r+1$ edge inequalities gives twice the cycle sum at least $2r+1$, so that integer sum is at least $r+1$. A minimum of the positive-cost linear objective over $P(G)$ exists in a bounded sublevel set and has an optimal vertex, by the elementary extreme-point theorem for linear programming. That vertex therefore cannot be integral.
\end{proof}
The result is the familiar graph-cover characterization, not the sought classification of all $0$--$1$ matrices. In particular, total unimodularity is not necessary for ideality once redundant rows are permitted. Take the three identity rows and the three pair rows $(1,1,0),(1,0,1),(0,1,1)$. Its covering polyhedron is exactly $\{x:x_i\ge1\}$, whose only vertex is integral. Yet the determinant of the submatrix formed by the pair rows has absolute value 2, so the whole matrix is not totally unimodular.
\subsection{Verification and remaining gaps}
The triangle has fractional cover value $3/2$ and integer value 2; the redundant-row example has its unique vertex at $(1,1,1)$ and a determinant of absolute value 2. These values were checked exactly. The attempted equivalence ``ideal if and only if totally unimodular'' fails on that example. The perturbation proof uses bipartite tight-edge components and does not extend directly to rows with three or more ones; the general ideal/minimally nonideal structure remains unresolved by this attempt.

\section{References and source attribution}
Primary workshop report: \url{https://doi.org/10.4171/owr/2005/50}. Relevant passages read: 2005/50 pp. 2876--2878, ideal covering matrices, minimally nonideal matrices and structural results already reported.

All retrievals were checked on 11 October 2026. Original question attribution remains in the preserved Sources section. External theorems are attributed and not claimed as new discoveries or fully verified merely by citation.


\paragraph{Integration review.} Codex (GPT-6; exact serving revision not exposed) reviewed the written arguments and expanded layout and source-reference wrappers. Repository metadata and the production web renderer were checked using Python 3.13 and JavaScript with MathJax 3.2.2. This model-assisted review is not independent human proof verification; no PDF compilation or formal proof checking is asserted.

\section{Completion Estimate}
\noindent\textbf{COMPLETION: 0\%}
No substantiated new progress toward the entire original target is claimed. The worked results and failed route are delimited above; a known result, illustrative model, or auxiliary obstruction is not counted as original-target completion.

\section{Final status and attribution}
\textbf{Attempt status: unresolved.} This describes this research attempt, not a new ruling on the source catalogue's status or on external literature claims.
The mathematics and this self-audit were generated by OpenAI \texttt{gpt-6-astra}, reasoning \texttt{ultra}, on 11 October 2026. Codex (GPT-6; exact serving revision not exposed) orchestrated the batch and contributed formatting/repository review. No human mathematical authorship, independent human proof checking, or formal verification is claimed. The historical preparation paragraphs in the reproduced catalogue body concern that body only.

Contribution prompt A was applied to the complete definition and supplied source record, followed by the adversarial checks of prompt B. Local reproduction files are \texttt{batch016/checks.py}, \texttt{computations.json}, and the per-ID audit JSON. They were not public links at preparation time. Numerical checks are finite and do not prove unrestricted results. No LaTeX engine was found locally; compilation was not performed by this research agent.

\end{document}
